\section{Conclusion}
\label{sec:conclusion}

\subsection{Summary}

We have characterised the 2030 reapportionment shock for the \ar{}
bisection algorithm.
Using Census Bureau projections and simulation on 2020 census-tract graphs,
we find:

\begin{enumerate}
\item \textbf{Texas is the most disruptive case.}
      $38 = 2 \times 19 \to 41$ (prime) changes the tree from a 2-level
      to a flat structure.
      Simulated Hamming distance: $d_{\rm Ham} = 0.23$.

\item \textbf{Florida is also dramatically affected.}
      $28 = 4 \times 7 \to 29$ (prime) produces a similar
      composite-to-prime transition.
      Simulated $d_{\rm Ham} = 0.19$.

\item \textbf{Reapportionment disruption is less than MCMC mixing.}
      The worst-case reapportionment disruption ($d_{\rm Ham} = 0.23$)
      is equivalent to approximately 10 \recom{} steps for Texas.
      A fully-mixed ensemble (10{,}000+ steps) produces $d_{\rm Ham} \approx 1$.

\item \textbf{Political redistricting is more disruptive.}
      Contested political redistricting cycles change 35--55\% of tract
      assignments, $1.5$--$2.4\times$ more than the worst algorithmic
      reapportionment.

\item \textbf{Fresh recomputation is the correct statutory design.}
      The DIA's decadal recomputation produces maps optimal for the current
      apportionment and geography.
      No continuity bonus or grandfathering clause is needed.
\end{enumerate}

\subsection{The Arithmetic of Disruption}

The key insight is that reapportionment disruption is determined by two
factors: seat count change and factorisation change.
A state gaining one seat from a composite $k$ to another composite $k+1$
(e.g., Georgia $14 \to 15$, both composite) has disruption $d_{\rm Ham} \approx 0.08$.
A state moving from composite to prime (Texas, Florida) has disruption
$d_{\rm Ham} \approx 0.20$--$0.23$.

This arithmetic is predictable and transparent.
States and administrators can compute the expected disruption years in advance
from Census Bureau projections.

\subsection{Open Questions}

\begin{enumerate}
\item \textbf{2030 census data.}
      The simulations use 2020 data with new seat counts.
      The actual 2030 disruption will be larger due to population redistribution.
      A 2030 simulation requires the 2030 census, expected in 2031.

\item \textbf{Optimal transition strategies.}
      If continuity is valued, is there an \ar{} variant that
      would minimise $d_{\rm Ham}$ between the old and new maps
      subject to compactness?
      This is an open algorithmic question.

\item \textbf{The prime frequency problem.}
      Among the 50-state seat counts, primes appear frequently
      (current primes: $\{7, 11, 13, 17, 19, 23, 29, 31, 37, 41\}$).
      A theory of reapportionment stability that accounts for prime
      density in the relevant range $[1, 60]$ would be useful.
      The prime number theorem gives density $\sim 1/\ln(k)$, so
      roughly $1/\ln(40) \approx 27\%$ of seat counts near 40 are prime.
\end{enumerate}

\subsection{Practical Advice}

For states expecting seat changes in 2030, we recommend:
\begin{enumerate}
\item Pre-compute the 2030 bisection tree from current projections.
\item Run a preliminary simulation using 2020 data with the projected
      seat count to estimate boundary disruption.
\item Publish the projected changes for public comment before the 2030
      redistricting cycle begins.
\end{enumerate}

This transparency is consistent with the DIA's overall commitment to
deterministic, auditable redistricting.
